\section{Method}
\label{sec:method}

\subsection{Control task}
Each episode is one boom move: the boom starts at rest at $\theta_0$ and must reach a
target angle $\theta^*$ within 10\,s, with $\theta_0$ and $\theta^*$ drawn uniformly over the
end-stop range (with a 0.05\,rad margin). The controller acts every $\Delta t=50$\,ms with a
normalised command $u\in[-1,1]$ (Section~\ref{sec:model}) and observes
\begin{equation}
o_k = \big[\,\theta,\ \dot\theta,\ e,\ \textstyle\int e\,dt\,\big],\qquad e=\theta^*-\theta,
\label{eq:obs}
\end{equation}
with the error integral clipped to $\pm1$ and reset each episode.
% Labs 03-07 used [theta, omega, e] and theta0 in [-0.4, 0.9], target in [-0.3, 0.9] rad.

\subsection{Reward and energy weight}
The reward trades accuracy against actuator power,
\begin{equation}
r_k = -\Big(e^2 + 0.01\,\dot\theta^{\,2} + w\,\frac{P_k}{P_\mathrm{ref}}\Big),
\label{eq:reward}
\end{equation}
where $P_k=|\tau_\mathrm{cyl}\dot\theta|$ is the mechanical power over the step and $w\ge0$ the
energy weight. $P_\mathrm{ref}$ normalises the power: $250$\,kW on the placeholder plant and the
pump corner power $Q\,p_\mathrm{relief}=271.7$\,kW on the real-parameter plant.
The first experiments penalised the squared command $w\,u^2$ instead of the power.
The operating point $w=0.1$ was chosen from a sweep of ten weights between 0 and 2
(three seeds each): it reduces the energy by 15\,\% relative to $w=0$
(131.6 vs 155.4\,kJ) at a similar mean error ($0.58^\circ$ vs $0.67^\circ$);
above $w\approx0.5$ the mean error rises past $2^\circ$.
% Source: README Lab 04, rl/04_knee_sweep.csv (u^2 penalty, placeholder plant, 20k steps).
% TODO: w keeps its value when the reward switches from u^2 to power (08b); say so in the text.

\subsection{Soft Actor-Critic}
The policy is trained with Soft Actor-Critic (SAC) \cite{haarnoja2018sac}, an off-policy
maximum-entropy actor-critic method for continuous actions, using the stable-baselines3
implementation \cite{raffin2021sb3} with its default hyperparameters and a multilayer-perceptron
policy. Each model is trained for $10^5$ steps (500 episodes); every configuration is trained
with three random seeds.
% TODO: list the SB3 2.9.0 defaults used; add both references to refs.bib.

\subsection{Plain and residual SAC}
Two policy structures are compared.
In \emph{plain} SAC the policy output is the command, $u=a$.
In \emph{residual} SAC the policy corrects a model-based base controller,
\begin{equation}
u = \mathrm{clip}\big(u_\mathrm{base} + a,\,-1,\,1\big),
\label{eq:residual}
\end{equation}
so that $a=0$ reproduces the base controller exactly.
On the placeholder plant the base controller is a PID with gravity feed-forward computed from
the nominal mass ($K_p=2$, $K_i=0.1$, $K_d=1$, integral frozen at saturation).
On the real-parameter plant it is the PID baseline of Section~\ref{sec:baselines} at its
balanced operating point.

\subsection{Domain randomisation}
To make the policy robust to model error, the plant is randomised at every reset.
On the placeholder plant: the boom mass is scaled by $\mathcal{U}(0.4,1.3)$; with probability 0.5
the torque limit becomes angle dependent (proportional to the cylinder moment arm); with
probability 0.3 a payload step multiplies the mass by $\mathcal{U}(0.9,1.3)$ at a random step
between 50 and 150.
On the real-parameter plant: the bucket tip mass is drawn from $\mathcal{U}(1120,3406)$\,kg
(empty to full bucket); with probability 0.3 it jumps to a new value from the same range at a
random step between 50 and 150; with probability 0.5 the torque limit is scaled by
$\mathcal{U}(0.7,1.0)$; and the inertia is scaled by $\mathcal{U}(0.38,1.5)$.
In a second variant the measured angle and speed are also corrupted with Gaussian noise whose
level is drawn per episode from $\{0, 0.2, 0.5, 1.0\}^\circ$; the reward keeps the true state.
% Empty bucket 1120 kg OFFICIAL, full 3406 kg ESTIMATED (1.27 m^3 x 1800 kg/m^3).

\subsection{PID baselines}
\label{sec:baselines}
On the placeholder plant the baseline is a PID with gravity feed-forward,
$u=\mathrm{clip}(K_p e + K_i\!\int\! e\,dt - K_d\dot\theta + u_\mathrm{ff})$, with gains from a grid
search (lowest mean energy among gain sets whose worst error stays below 20\,cm).
On the real-parameter plant the baseline is a gain-scheduled PID whose feed-forward uses the
real gravity model (\ref{eq:grav}) with an online estimate $\hat m$ of the bucket tip mass,
\begin{equation}
\hat m_\mathrm{raw} = \frac{\tau_\mathrm{cyl} - J\dot\omega - \tau_f - \tau_g(\theta,0)}{g\,x_\mathrm{tip}(\theta)},
\label{eq:mhat}
\end{equation}
filtered by an adaptive low-pass filter that restarts when a change in mass is detected.
Three operating points were selected from the energy--move-time Pareto front:
fast ($K_p=5$, $K_d=2$), balanced ($K_p=1.5$, $K_d=1$) and economical ($K_p=0.8$, $K_d=0.8$).
All controller tuning used a fixed set of 20 scenarios that is not used for evaluation.
% Estimator: adapt_Tf1_c1_n2 (T_f = 1 s, c = 1, n_hold = 2), Lab 11c, tuning set only.
